\ssr{ЗАКЛЮЧЕНИЕ}

В ходе выполнения лабораторной работы были реализованны и исследованы алгоритмы умножения матриц. Все цели лабораторной работы были достигнуты, а задачи выполнены.
\begin{enumerate}
	\item описаны алгоритмы умножения матриц: стандартный и Винограда.
	\item представлены блок-схемы трёх рассматриваемых алгоритмов: стандартного, классического Винограда и его оптимизированной версии.
	\item введена модель вычислений и выполнена оценка трудоёмкости алгоритмов на основе их схем или реализаций.
	\item реализован каждый алгоритм.
	\item проведены замеры процессорного времени выполнения алгоритмов.
	\item выполнены сравнительный анализ алгоритмов на основе данных замеров.
\end{enumerate}



Такое преимущество достигается благодаря двум ключевым улучшениям:
\begin{enumerate}
	\item объединение циклов --- исключены избыточные проходы по данным, что снижает накладные расходы на управление итерациями.
	\item замена операций умножения на битовые сдвиги --- операция сдвига (\lstinline{<< 1}) выполняется процессором значительно быстрее, чем умножение на 2.
\end{enumerate}



Наилучшие результаты продемонстрировал оптимизированный алгоритм Винограда: он оказался быстрее классического варианта примерно в 1,5 раза и превосходит базовую версию алгоритма Винограда почти в 1,2 раза.
